\section{Potencias con exponentes racionales}

Hasta ahora hemos dado significado a expresiones como
\[
a^3,\qquad a^0,\qquad a^{-2},
\]
es decir, a potencias cuyos exponentes son números enteros.

Queremos ahora comprender qué podría significar que el exponente sea una fracción:
\[
a^{1/2},
\qquad
a^{1/3},
\qquad
a^{2/5}.
\]

La idea será la misma que hemos seguido hasta ahora: no inventaremos una regla nueva de manera aislada, sino que buscaremos una definición compatible con las propiedades de las potencias que ya conocemos.

Para evitar dificultades relacionadas con los signos, comenzaremos trabajando con
\[
a>0.
\]

\subsection{El significado de exponentes fraccionarios}

Comencemos con un exponente sencillo:
\[
a^{1/2}.
\]

Sabemos que una potencia de otra potencia satisface
\[
(a^m)^n=a^{mn}.
\]

Si queremos que esta propiedad siga siendo válida al introducir exponentes fraccionarios, entonces debería ocurrir que
\[
\left(a^{1/2}\right)^2
=
a^{(1/2)\cdot2}
=
a^1
=
a.
\]

Por lo tanto, \(a^{1/2}\) debe representar un número cuyo cuadrado sea \(a\).

Pero ya tenemos una operación que responde exactamente a esa pregunta:
\[
\sqrt a.
\]

Por eso definimos
\[
a^{1/2}=\sqrt a.
\]

El mismo razonamiento puede hacerse con cualquier entero positivo \(q\). Si queremos que
\[
\left(a^{1/q}\right)^q
=
a^{(1/q)q}
=
a,
\]
entonces \(a^{1/q}\) debe ser el número cuya potencia \(q\)-ésima es \(a\).

Por lo tanto,
\[
\boxed{a^{1/q}=\sqrt[q]{a}}.
\]

Así, por ejemplo,
\[
16^{1/2}=\sqrt{16}=4
\]
y
\[
8^{1/3}=\sqrt[3]{8}=2.
\]

El exponente fraccionario \(1/q\) no introduce entonces una operación completamente nueva: es otra manera de escribir una raíz.

\subsubsection{Exponentes fraccionarios de la forma \(p/q\)}

Consideremos ahora un exponente como
\[
\frac{p}{q},
\]
donde inicialmente supondremos que \(p\) y \(q\) son naturales positivos.

Podemos escribir
\[
\frac pq=p\cdot\frac1q.
\]

Utilizando nuevamente la propiedad de una potencia elevada a otra potencia,
\[
a^{p/q}
=
a^{p(1/q)}
=
\left(a^{1/q}\right)^p.
\]

Como acabamos de establecer que
\[
a^{1/q}=\sqrt[q]{a},
\]
obtenemos
\[
\boxed{
a^{p/q}
=
\left(\sqrt[q]{a}\right)^p
}.
\]

Esto nos da una forma concreta de interpretar el exponente racional:

\begin{enumerate}
    \item calculamos la raíz \(q\)-ésima de la base;
    \item elevamos el resultado a la potencia \(p\).
\end{enumerate}

Por ejemplo,
\[
8^{2/3}
=
\left(\sqrt[3]{8}\right)^2
=
2^2
=
4.
\]

\subsection{¿Podemos realizar las operaciones en el orden contrario?}

Ya vimos que
\[
a^{p/q}
=
\left(\sqrt[q]{a}\right)^p.
\]

Esto significa que podemos calcular primero la raíz \(q\)-ésima y después elevar el resultado a la potencia \(p\).

Pero observemos nuevamente el exponente:
\[
\frac pq
=
p\cdot\frac1q.
\]

Como la multiplicación es conmutativa, también podemos escribir
\[
\frac pq
=
\frac1q\cdot p.
\]

Los dos factores son exactamente los mismos. Solo hemos cambiado su orden.

Y aquí ocurre algo interesante.

Si una potencia elevada a otra potencia multiplica sus exponentes, entonces
\[
a^{p/q}
=
a^{(1/q)p}
=
\left(a^p\right)^{1/q}.
\]

Pero ya sabemos que un exponente de la forma \(1/q\) representa una raíz \(q\)-ésima:
\[
x^{1/q}=\sqrt[q]{x}.
\]

Por lo tanto,
\[
\left(a^p\right)^{1/q}
=
\sqrt[q]{a^p}.
\]

Así obtenemos
\[
\boxed{
a^{p/q}
=
\left(\sqrt[q]{a}\right)^p
=
\sqrt[q]{a^p}
}.
\]

Podemos verlo directamente en los exponentes:
\[
p\cdot\frac1q
=
\frac1q\cdot p.
\]

En el primer orden,
\[
\left(a^{1/q}\right)^p,
\]
calculamos primero la raíz y después la potencia.

En el segundo,
\[
\left(a^p\right)^{1/q},
\]
calculamos primero la potencia y después la raíz.

No son dos reglas diferentes. Son las mismas operaciones realizadas en distinto orden.

Veamos un ejemplo:
\[
27^{2/3}.
\]

Podemos calcular primero la raíz:
\[
27^{2/3}
=
\left(\sqrt[3]{27}\right)^2
=
3^2
=
9.
\]

O podemos calcular primero la potencia:
\[
27^{2/3}
=
\sqrt[3]{27^2}
=
\sqrt[3]{729}
=
9.
\]

En ambos casos llegamos al mismo resultado.

Normalmente convendrá elegir el camino que produzca cálculos más sencillos. En este ejemplo,
\[
\sqrt[3]{27}=3
\]
hace que calcular primero la raíz sea mucho más cómodo que comenzar con
\[
27^2=729.
\]

\subsection{La fracción utilizada no puede cambiar el resultado}

Hay todavía un detalle importante.

Sabemos que una misma fracción puede escribirse de muchas maneras:
\[
\frac12=\frac24=\frac36=\cdots
\]

Por lo tanto, si utilizamos estas fracciones como exponentes, todas deben producir el mismo resultado.

Por ejemplo,
\[
16^{1/2}=4.
\]

Pero también podemos escribir
\[
\frac12=\frac24.
\]

Entonces deberíamos obtener igualmente
\[
16^{2/4}=4.
\]

Y, en efecto,
\[
16^{2/4}
=
\sqrt[4]{16^2}
=
\sqrt[4]{256}
=
4.
\]

No parece haber ningún problema.

Pero queremos asegurarnos de que esto ocurra siempre.

Supongamos que tenemos el exponente
\[
\frac pq
\]
y multiplicamos su numerador y su denominador por el mismo número natural positivo \(k\):
\[
\frac pq=\frac{pk}{qk}.
\]

Utilizando la primera fracción obtenemos
\[
a^{p/q}
=
\sqrt[q]{a^p}.
\]

Ahora preguntémonos: ¿qué sucede si usamos la segunda?
Obtendríamos
\[
a^{pk/qk}
=
\sqrt[qk]{a^{pk}}.
\]

A primera vista las dos raíces parecen bastante diferentes:
\[
\sqrt[q]{a^p}
\qquad\text{y}\qquad
\sqrt[qk]{a^{pk}}.
\]

Para comprobar si son iguales, planteamos la igualdad:
\begin{align*}
  \sqrt[q]{a^p} &\stackrel{?}{=} \sqrt[qk]{a^{pk}} \\ 
  \left(\sqrt[q]{a^p}\right)^{qk}
  &\stackrel{?}{=}
  \left(\sqrt[qk]{a^{pk}}\right)^{qk}
\end{align*}

Aplicando las propiedades de los exponentes y las raíces que hemos visto:
\begin{align*}
  \left(\sqrt[q]{a^p}\right)^{qk}
  &\stackrel{?}{=}
  \left(a^{pk}\right)^{\frac{qk}{qk}} \\ 
  \left(a^p\right)^{\frac{qk}{q}}
  &\stackrel{?}{=}
  a^{pk} \\ 
  \left(a^p\right)^k
  &\stackrel{?}{=}
  a^{pk} \\ 
  a^{pk}
  &\stackrel{?}{=}
  a^{pk}.
\end{align*}

La última igualdad es verdadera. Además, como estamos trabajando con raíces principales, las dos expresiones originales son positivas. Por lo tanto, si sus potencias \(qk\)-ésimas son iguales, ellas también deben ser iguales.

Así,
\[
\sqrt[q]{a^p}
=
\sqrt[qk]{a^{pk}}.
\]

Y, en consecuencia,
\[
\boxed{
a^{p/q}=a^{pk/qk}
}.
\]

Así que ampliar una fracción no cambia el valor de la potencia.

Por ejemplo,
\[
\frac12=\frac24,
\]
y por eso
\[
16^{1/2}
=
16^{2/4}
=
4.
\]

Lo mismo ocurriría con
\[
\frac12=\frac36,
\qquad
\frac23=\frac46,
\qquad
\frac34=\frac68,
\]
o con cualquier otra fracción equivalente.

Esto es fundamental: el valor de
\[
a^{p/q}
\]
depende del número racional \(\frac pq\), no de la manera particular en que hayamos escrito esa fracción.

\subsection{Numerador cero o negativo}

Hasta ahora supusimos que \(p>0\). Podemos completar la definición utilizando lo que ya sabemos sobre exponentes enteros.

Si
\[
p=0,
\]
entonces
\[
\frac pq=0,
\]
y recuperamos
\[
a^{0/q}=a^0=1.
\]

Si \(p<0\), utilizamos la misma idea que empleamos para los exponentes enteros negativos: el exponente negativo indica el inverso multiplicativo.

Así,
\[
a^{-p/q}
=
\frac{1}{a^{p/q}}.
\]

Por ejemplo,
\[
8^{-2/3}
=
\frac{1}{8^{2/3}}.
\]

Como
\[
8^{2/3}
=
\left(\sqrt[3]{8}\right)^2
=
2^2
=
4,
\]
obtenemos
\[
8^{-2/3}
=
\frac14.
\]

De este modo, para una base positiva \(a\), los exponentes racionales quedan conectados con operaciones que ya conocemos: potencias, raíces e inversos multiplicativos.

\subsection{Base cero}

Para \(r\in\mathbb Q\) positivo, podemos definir
\[
0^r=0.
\]
Pero \(0^0\) queda sin definir en este capítulo y \(0^r\) tampoco está definido cuando \(r<0\), porque un exponente negativo exigiría dividir por una potencia de cero.

\subsection{Bases negativas}

Hasta ahora hemos trabajado principalmente con bases positivas. Con ellas, las raíces y los exponentes racionales se comportan de una manera bastante cómoda.

Con las bases negativas debemos prestar un poco más de atención.

Comencemos con algo que ya conocemos:
\[
(-8)^{1/3}.
\]

Como el exponente \(1/3\) indica una raíz cúbica,
\[
(-8)^{1/3}
=
\sqrt[3]{-8}
=
-2.
\]

No hay ningún problema: una cantidad negativa puede tener una raíz de índice impar, ya que
\[
(-2)^3=-8.
\]

Algo parecido ocurre con
\[
(-8)^{2/3}.
\]

Primero calculamos la raíz cúbica y después elevamos al cuadrado:
\[
(-8)^{2/3}
=
\left(\sqrt[3]{-8}\right)^2
=
(-2)^2
=
4.
\]

Observemos algo interesante: la base sigue siendo negativa, pero ahora el resultado es positivo.

Esto no es extraño. La raíz cúbica nos dio \(-2\), y al elevarlo al cuadrado desapareció el signo negativo.

En general, si escribimos el exponente racional en su forma irreducible
\[
\frac pq,
\]
el denominador \(q\) nos indica qué raíz debemos calcular.

Si \(q\) es impar, podemos calcular una raíz \(q\)-ésima de un número negativo dentro de los números reales.

Por ejemplo,
\[
\sqrt[3]{-8}=-2,
\qquad
\sqrt[5]{-32}=-2.
\]

Después, el numerador \(p\) nos indica a qué potencia debemos elevar ese resultado. Según sea esa potencia par o impar, el signo final puede cambiar.

Pero si \(q\) es par, aparece un problema.

Por ejemplo,
\[
(-16)^{1/4}
=
\sqrt[4]{-16}.
\]

No existe ningún número real cuya cuarta potencia sea \(-16\), porque una potencia par de un número real nunca es negativa.

Por lo tanto,
\[
(-16)^{1/4}
\]
no representa un número real.

\subsubsection*{Una pequeña trampa}

Aquí aparece una razón muy importante para reducir siempre la fracción del exponente.

Sabemos que
\[
\frac13=\frac26.
\]

Por lo tanto,
\[
(-8)^{1/3}
\qquad\text{y}\qquad
(-8)^{2/6}
\]
tienen exactamente el mismo exponente.

Ya sabemos que
\[
(-8)^{1/3}=-2.
\]

Pero imaginemos que tomamos \(2/6\) sin reducirlo y aplicamos mecánicamente la expresión
\[
a^{p/q}=\sqrt[q]{a^p}.
\]

Obtendríamos
\[
(-8)^{2/6}
\stackrel{?}{=}
\sqrt[6]{(-8)^2}
=
\sqrt[6]{64}
=
2.
\]

¡Pero acabamos de obtener \(2\) cuando antes habíamos obtenido \(-2\)!

El problema no está en que
\[
\frac13
\quad\text{y}\quad
\frac26
\]
sean exponentes diferentes. Son exactamente el mismo número racional.

El problema es que, al utilizar \(2/6\) sin reducir, introdujimos artificialmente una raíz de índice par. Al elevar primero \(-8\) al cuadrado, además, perdimos el signo negativo antes de calcular la raíz.

Por eso, cuando la base es negativa, primero debemos escribir el exponente racional en su forma irreducible:
\[
\frac26=\frac13.
\]

Entonces calculamos
\[
(-8)^{2/6}
=
(-8)^{1/3}
=
\sqrt[3]{-8}
=
-2.
\]

Así evitamos que una forma innecesariamente complicada de escribir la fracción nos lleve por un camino equivocado.

Podemos encontrar una situación todavía más peligrosa.

Consideremos
\[
(-5)^{2/4}.
\]

Si observamos solamente la expresión tal como está escrita, podríamos pensar en calcular primero la potencia:
\[
(-5)^{2/4}
\stackrel{?}{=}
\sqrt[4]{(-5)^2}.
\]

Entonces,
\[
\sqrt[4]{(-5)^2}
=
\sqrt[4]{25}
=
\sqrt{5}.
\]

¡Parece que hemos encontrado un resultado perfectamente válido! Pero algo huele mal de este resultado. Observemos que ocurre si intentamos verificar la solución. 

Supuestamente decimos que 
\[
  (-5)^{2/4} = \sqrt5
\]
En otras palabras, aplicando la definición de raíz:
\[
  (-5)^2 = \left(\sqrt{5}\right)^4
\]
Es decir, 
\[
  (-5)^2 = (5)^2
\]
Lo cual es verdad. Sin embargo observemos un detalle sutil: los números de la base no son iguales, a pesar de que el resultado final si lo es.

Pero hagamos algo que siempre debemos hacer con un exponente racional: reduzcamos primero la fracción.

Como
\[
\frac24=\frac12,
\]
la expresión original es en realidad
\[
(-5)^{1/2}.
\]

Y eso significa
\[
\sqrt{-5},
\]
que no es un número real.

Así que hemos llegado a una situación bastante grave. Por un camino logramos encontrar que:
\[
(-5)^{2/4}
\stackrel{?}{=}
\sqrt5
\]
mientras que por otro camino llegamos a:
\[
(-5)^{2/4}
=
(-5)^{1/2}
\]
que nos muestra que la expresión no representa ningún número real.

¿Qué ocurrió?

Al elevar primero \(-5\) al cuadrado obtuvimos
\[
(-5)^2=25.
\]

En ese instante desapareció el signo negativo. Después calculamos una raíz cuarta de un número positivo y todo pareció funcionar.

Pero habíamos ocultado el verdadero problema.

El exponente
\[
\frac24
\]
no es un número racional diferente de
\[
\frac12.
\]

Son exactamente el mismo número:
\[
\frac24=\frac12.
\]

Por lo tanto, no podemos permitir que una forma distinta de escribir el mismo exponente cambie si la potencia existe o incluso produzca un resultado diferente.

Este ejemplo nos muestra por qué, con bases negativas, reducir primero el exponente racional no es solamente una comodidad:
\[
\boxed{\text{es necesario para interpretar correctamente la potencia.}}
\]

Así, ante una expresión como
\[
(-5)^{2/4},
\]
primero escribimos
\[
\frac24=\frac12.
\]

Entonces vemos inmediatamente que
\[
(-5)^{2/4}
=
(-5)^{1/2}
\]
no representa un número real.

El valor \(\sqrt5\) que apareció antes no era una segunda respuesta posible. Fue el resultado de seguir un camino algebraico que, en este caso, no era válido.

Podemos resumir la idea de esta manera:

Si la base es negativa, primero reducimos el exponente
\[
\frac pq
\]
a su forma irreducible.

Luego observamos su denominador:
\[
\boxed{
\begin{array}{c}
q\text{ impar} \quad\longrightarrow\quad
\text{la potencia puede representar un número real},\\[4pt]
q\text{ par} \quad\longrightarrow\quad
\text{la potencia no representa un número real}.
\end{array}
}
\]
Por esta razón, cuando estudiemos las propiedades generales de los exponentes racionales utilizaremos normalmente bases positivas. Así podremos concentrarnos en las propiedades mismas sin detenernos constantemente a comprobar qué raíces existen en los números reales.
